\section{记忆的层级与流转}\label{sec:tiers}

\subsection{三层记忆如何协作}

助手同时维护三层记忆，它们各司其职。\textbf{短期记忆}是当前对话的上下文窗口，容量
有限、随会话结束。\textbf{中期记忆}在上下文即将溢出时接管：对话中确立的任务目标、
做出的决策、改动过的文件、遇到的错误与解决办法，会在压缩时被抽取为结构化记录，
跨会话保留。\textbf{长期记忆}则是前面章节介绍的偏好库、知识库与语义索引，永久保存
并持续注入后续对话。

日常使用中你不需要干预这套流转，但理解它有助于解释两个常见现象。其一，「助手好像
忘了很早之前的对话细节」——细节可能还保留在中期记忆中，只是不在当前窗口里，你可以
让助手「查一下之前关于 xx 的记录」主动召回。其二，「压缩之后助手仍记得我决定过
什么」——这正是压缩退役网关在起作用，重要决策以高于普通内容的重要性被保留。

\subsection{中期记忆的读取与晋升}

你可以通过自然语言查询中期记忆，例如「这个会话里我们做过哪些决定」「之前处理 xx
错误的方法是什么」。中期记忆中反复被访问、重要性高的记录，最终应晋升为长期记忆；
该晋升机制与淘汰策略的实现进度见 \todo{中期记忆晋升机制落地后更新本节操作说明}。
